\documentclass[11pt,fontset=none]{ctexart}
\usepackage[a4paper,margin=0.8in]{geometry}
\usepackage{longtable}
\usepackage{booktabs}
\usepackage{array}
\usepackage{hyperref}
\usepackage{xcolor}
\usepackage{fancyhdr}
\setlength{\headheight}{15pt}
\setCJKmainfont{STSong}
\setCJKsansfont{STSong}
\setCJKmonofont{STSong}
\pagestyle{fancy}
\fancyhf{}
\lhead{Cell Biology Glossary Table}
\rhead{notebook + PDFs}
\cfoot{\thepage}

\newcolumntype{L}[1]{>{\raggedright\arraybackslash}p{#1}}

\title{Cell Biology 资料中出现的蛋白质、术语与方法简表}
\author{}
\date{2026-04-27}

\begin{document}
\maketitle

\section*{说明}
本表依据当前文件夹中的 \texttt{notebook.md}、\texttt{CB-Ch4-041626.pdf} 与 \texttt{CB-Ch67-042326.pdf} 整理。目标是尽量覆盖资料中\textbf{明确提到}的蛋白质、重要术语和方法名称；去掉了大量纯小分子、普通细胞器名称和重复项。功能描述故意保持非常简短，方便临考速查。

\small

\section*{1. 细胞骨架与运动}
\begin{longtable}{L{3.2cm} L{3.8cm} L{8.5cm}}
\toprule
\textbf{名称} & \textbf{中文/类别} & \textbf{功能（极简）} \\
\midrule
\endfirsthead
\toprule
\textbf{名称} & \textbf{中文/类别} & \textbf{功能（极简）} \\
\midrule
\endhead
\bottomrule
\endfoot
Actin & 肌动蛋白 & 形成微丝，参与细胞形态、迁移、收缩和胞质分裂。 \\
$\alpha$-actin & $\alpha$-肌动蛋白 & 主要在肌细胞中，负责收缩。 \\
$\beta$-actin & $\beta$-肌动蛋白 & 常见于非肌细胞，偏向迁移与前沿延伸。 \\
$\gamma$-actin & $\gamma$-肌动蛋白 & 常见于非肌细胞，偏向皮层网络稳定与形态维持。 \\
Arp2/3 complex & Arp2/3复合体 & 成核分支状 actin network。 \\
Formin & 成束相关成核蛋白 & 成核并延长长而直的 actin filament。 \\
Profilin & 辅助蛋白 & 促进 actin 单体加到正端。 \\
Thymosin & 胸腺素样 actin 结合蛋白 & 结合 actin 单体，限制其聚合。 \\
CapZ & 正端帽蛋白 & 封住 actin filament 的 plus end。 \\
Cofilin & 辅助蛋白 & 扭曲并削弱 actin filament，促进切断和去聚合。 \\
NPF & 成核促进因子 & 激活 Arp2/3 complex。 \\
Fimbrin & 交联蛋白 & 形成紧密平行 actin bundles。 \\
$\alpha$-actinin & 交联蛋白 & 形成较疏松的 contractile actin bundles。 \\
Filamin A & 交联蛋白 & 交联 actin network；突变与神经迁移异常相关。 \\
Myosin II & 肌球蛋白II & actin 马达蛋白；驱动肌肉收缩和 contractile bundle 收缩。 \\
Myosin V & 肌球蛋白V & 沿 actin 运输细胞器/货物。 \\
Tropomyosin & 原肌球蛋白 & 沿薄丝排列，调节 myosin 对 actin 的结合。 \\
Troponin T & 肌钙蛋白T & 结合 tropomyosin。 \\
Troponin I & 肌钙蛋白I & 抑制 actin-myosin 相互作用。 \\
Troponin C & 肌钙蛋白C & 结合 Ca$^{2+}$，启动骨骼肌收缩调控。 \\
Titin & Titin/连接蛋白 & 肌节分子弹簧，帮助定位粗丝并提供弹性。 \\
Nebulin & Nebulin & 参与薄丝长度控制。 \\
Tropomodulin & Tmod/减端帽蛋白 & 封住薄丝 minus end，稳定长度。 \\
Calmodulin & 钙调蛋白 & 结合 Ca$^{2+}$；平滑肌中激活 MLCK 通路。 \\
FtsZ & 原核 tubulin 同源蛋白 & 形成 Z-ring，参与细菌胞质分裂。 \\
MreB & 原核 actin 同源蛋白 & 维持细菌细胞形状。 \\
ParM & 原核 actin 同源蛋白 & 参与质粒/染色体分离。 \\
Crescentin & 原核中间丝样蛋白 & 帮助维持弯曲细菌形态。 \\
Tubulin & 微管蛋白 & 构成微管。 \\
$\alpha$-tubulin & $\alpha$-微管蛋白 & 构成微管异二聚体；其 GTP 不交换。 \\
$\beta$-tubulin & $\beta$-微管蛋白 & 构成微管异二聚体；其 GTP 水解驱动动态不稳定。 \\
$\gamma$-tubulin & $\gamma$-微管蛋白 & 微管成核。 \\
MAP2 & 微管相关蛋白2 & 稳定并组织微管。 \\
Tau & Tau蛋白 & 稳定并组织微管。 \\
Kinesin & 驱动蛋白 & 多数朝微管 plus end 运输。 \\
Kinesin-13 & 微管动力学马达 & 促进微管末端解聚。 \\
Dynein & 动力蛋白 & 多数朝微管 minus end 运输。 \\
Dynactin & Dynein 辅助复合体 & 连接 dynein 与膜性细胞器。 \\
Arp1 filament & Dynactin 组成部分 & 构成 dynactin 骨架的一部分。 \\
XMAP215 & 微管生长因子 & 把 tubulin 二聚体送到微管正端，促进生长。 \\
Katanin & 微管切割蛋白 & 切断微管。 \\
Lis1 & Dynein 结合蛋白 & 帮助核迁移；缺陷可致 lissencephaly。 \\
Rac & Rho 家族小GTP酶 & 促进 lamellipodia 和 branched actin network 形成。 \\
Rho & Rho 家族小GTP酶 & 促进 stress fiber 和 focal adhesion 形成。 \\
Keratins & 角蛋白家族 & 上皮细胞主要中间丝，提供机械强度。 \\
Lamins & 核纤层蛋白 & 构成 nuclear lamina。 \\
Lamin A & A型 lamin & 核纤层结构支持与染色质组织。 \\
Lamin C & A型 lamin & 核纤层结构支持与染色质组织。 \\
Lamin B1 & B型 lamin & 核纤层结构支持与核膜组织。 \\
Lamin B2 & B型 lamin & 核纤层结构支持与核膜组织。 \\
Plectin & Plakin 家族蛋白 & 连接 intermediate filaments 与微管/其他骨架。 \\
SUN proteins & SUN 蛋白 & 与 KASH 蛋白跨核膜连接核与细胞质骨架。 \\
KASH proteins & KASH 蛋白 & 与 SUN 蛋白一起形成核-细胞质力学连接。 \\
Ankyrin & 锚蛋白 & 把 spectrin 骨架连接到膜蛋白。 \\
Spectrin & 血影蛋白 & 皮层骨架成分，限制膜蛋白扩散并增强机械强度。 \\
ActA & 细菌表面蛋白 & 激活宿主 Arp2/3，劫持 actin 运动。 \\
\end{longtable}

\section*{2. 核、染色质与染色体}
\begin{longtable}{L{3.2cm} L{3.8cm} L{8.5cm}}
\toprule
\textbf{名称} & \textbf{中文/类别} & \textbf{功能（极简）} \\
\midrule
\endfirsthead
\toprule
\textbf{名称} & \textbf{中文/类别} & \textbf{功能（极简）} \\
\midrule
\endhead
\bottomrule
\endfoot
Histones & 组蛋白 & 与 DNA 形成核小体。 \\
Histone H1 & 连接组蛋白H1 & 结合 linker DNA，促进染色质进一步压缩。 \\
Histone H2A & 核心组蛋白H2A & 核小体八聚体成分。 \\
Histone H2B & 核心组蛋白H2B & 核小体八聚体成分。 \\
Histone H3 & 核心组蛋白H3 & 核小体八聚体成分；常带多种修饰。 \\
Histone H4 & 核心组蛋白H4 & 核小体八聚体成分。 \\
H3.3 & 组蛋白变体H3.3 & 参与表观遗传状态维持。 \\
CENP-A & 着丝粒组蛋白变体 & 定义 centromere 特化核小体。 \\
HP1 & 异染色质蛋白1 & 识别抑制性修饰，促进异染色质形成。 \\
RSC complex & 染色质重塑复合体 & ATP 依赖地滑动/重塑核小体。 \\
ING-PHD domain protein & 修饰读取蛋白 & 识别特定组蛋白甲基化状态。 \\
Reader proteins & 修饰读取蛋白 & 识别已有染色质标记。 \\
Writer proteins & 修饰写入蛋白 & 在邻近染色质上写入相同修饰。 \\
MyoD & 转录因子 & 与肌细胞命运和表观遗传记忆相关。 \\
\end{longtable}

\section*{3. 转录、RNA加工与基因调控}
\begin{longtable}{L{3.2cm} L{3.8cm} L{8.5cm}}
\toprule
\textbf{名称} & \textbf{中文/类别} & \textbf{功能（极简）} \\
\midrule
\endfirsthead
\toprule
\textbf{名称} & \textbf{中文/类别} & \textbf{功能（极简）} \\
\midrule
\endhead
\bottomrule
\endfoot
RNA polymerase & RNA聚合酶 & 以 DNA 为模板合成 RNA。 \\
RNA polymerase I & Pol I & 主要转录 rRNA。 \\
RNA polymerase II & Pol II & 主要转录 mRNA 和部分 noncoding RNA。 \\
RNA polymerase III & Pol III & 主要转录 tRNA、5S rRNA 和小 RNA。 \\
General transcription factors (TFII) & 通用转录因子 & 帮助 Pol II 在启动子起始转录。 \\
Sigma factor & $\sigma$ 因子 & 细菌中帮助 RNA polymerase 识别启动子。 \\
Topoisomerase & 拓扑异构酶 & 消除复制/转录产生的 superhelical tension。 \\
Phosphatase & 磷酸酶 & 5' capping 过程中处理 RNA 5' 端。 \\
Guanyl transferase & 鸟苷转移酶 & 给新生 RNA 加 5' cap。 \\
Methyltransferase & 甲基转移酶 & 对 cap 或染色质等进行甲基化。 \\
CPSF & 3'端加工因子 & 识别 poly(A) 信号并参与切割/加尾。 \\
CSTF & 3'端加工因子 & 参与 pre-mRNA 3' 端切割。 \\
ADAR & RNA 编辑酶 & 催化 A-to-I RNA editing。 \\
Dicer & RNA干扰相关酶 & 把双链 RNA 切成 siRNA/miRNA 前体。 \\
RISC & RNA诱导沉默复合体 & 利用小 RNA 引导抑制或降解靶 RNA。 \\
Tryptophan repressor & 色氨酸阻遏蛋白 & 细菌中抑制相关基因转录。 \\
CAP activator & CAP 激活蛋白 & 细菌中促进 lac operon 等转录激活。 \\
Bicoid & 转录调控蛋白 & 果蝇胚胎前后轴定位信息。 \\
Hunchback & 转录调控蛋白 & 参与果蝇胚胎图式建立。 \\
Kr\"uppel & 转录调控蛋白 & 参与果蝇胚胎条带调控。 \\
Giant & 转录调控蛋白 & 参与果蝇胚胎条带调控。 \\
Eve protein & Even-skipped 蛋白 & 果蝇胚胎分节调控。 \\
Eyeless & 转录调控蛋白 & 指定眼发育命运。 \\
Xist & lncRNA & 介导 X chromosome inactivation。 \\
\end{longtable}

\section*{4. 膜蛋白、通道与受体}
\begin{longtable}{L{3.2cm} L{3.8cm} L{8.5cm}}
\toprule
\textbf{名称} & \textbf{中文/类别} & \textbf{功能（极简）} \\
\midrule
\endfirsthead
\toprule
\textbf{名称} & \textbf{中文/类别} & \textbf{功能（极简）} \\
\midrule
\endhead
\bottomrule
\endfoot
Aquaporin & 水通道蛋白 & 高选择性转运水分子。 \\
Bacteriorhodopsin & 细菌视紫红质 & 光驱动质子泵。 \\
SGLT2 & Na$^+$-葡萄糖共转运蛋白2 & 肾脏重吸收葡萄糖；抗糖尿病药靶点。 \\
Serotonin transporter & 5-HT 转运体 & Na$^+$ 依赖再摄取 serotonin。 \\
Norepinephrine transporter & 去甲肾上腺素转运体 & Na$^+$ 依赖再摄取 norepinephrine。 \\
ABC transporters & ABC 转运蛋白家族 & ATP 驱动物质跨膜转运。 \\
Na$^+$/K$^+$ pump & 钠钾泵 & 建立 Na$^+$ / K$^+$ 梯度和膜电位。 \\
Ca$^{2+}$-ATPase & 钙泵 & 把 Ca$^{2+}$ 泵回 ER/SR 或细胞外。 \\
F$_1$F$_0$ ATP synthase & ATP 合酶 & 依靠质子梯度合成 ATP。 \\
Mechanosensitive channel & 机械敏感通道 & 保护细菌抵抗渗透压冲击。 \\
K$^+$ channel & 钾离子通道 & 选择性通透 K$^+$，参与静息电位。 \\
Voltage-gated Na$^+$ channel & 电压门控钠通道 & 介导动作电位快速去极化。 \\
Voltage-gated K$^+$ channel & 电压门控钾通道 & 参与复极化。 \\
Acetylcholine receptor & 乙酰胆碱受体 & 神经肌接头 ligand-gated cation channel。 \\
NMDA receptor & NMDA受体 & 允许 Ca$^{2+}$ 内流，参与 LTP。 \\
AMPA receptor & AMPA受体 & 介导快速兴奋性突触传递。 \\
GPI-anchored proteins & GPI 锚定蛋白 & 通过糖脂锚连接到膜外侧。 \\
\end{longtable}

\section*{5. 蛋白分选、细胞器导入与ER质量控制}
\begin{longtable}{L{3.2cm} L{3.8cm} L{8.5cm}}
\toprule
\textbf{名称} & \textbf{中文/类别} & \textbf{功能（极简）} \\
\midrule
\endfirsthead
\toprule
\textbf{名称} & \textbf{中文/类别} & \textbf{功能（极简）} \\
\midrule
\endhead
\bottomrule
\endfoot
Nucleoporins & 核孔蛋白 & 构成 nuclear pore complex。 \\
Importin & 核输入受体 & 识别 NLS cargo 并介导入核。 \\
Exportin & 核输出受体 & 识别 NES cargo 并介导出核。 \\
Ran & Ran GTPase & 为核进出口提供方向性。 \\
Ran-GAP & Ran GTP 酶激活蛋白 & 促进 Ran-GTP 水解。 \\
Ran-GEF & Ran 鸟苷交换因子 & 促进 Ran 结合 GTP。 \\
TOM complex & 线粒体外膜转位复合体 & 蛋白进入线粒体的总入口。 \\
TIM23 & 线粒体内膜转位复合体 & 导入基质蛋白或部分内膜蛋白。 \\
TIM22 & 线粒体内膜转位复合体 & 插入多跨膜载体蛋白到内膜。 \\
SAM complex & 外膜装配复合体 & 装配线粒体外膜 $\beta$-barrel 蛋白。 \\
OXA complex & 内膜插入复合体 & 将部分蛋白插入线粒体内膜。 \\
Hsp70 & Hsp70伴侣蛋白家族 & 帮助蛋白保持可转运状态或辅助导入/折叠。 \\
Cytosolic Hsp70 & 细胞质 Hsp70 & 保持线粒体前体蛋白未折叠，利于导入。 \\
Mitochondrial Hsp70 & 线粒体 Hsp70 & ATP 依赖地牵引蛋白进入基质。 \\
SRP & 信号识别颗粒 & 识别 ER signal peptide 并短暂停止翻译。 \\
SRP receptor & SRP受体 & 将携带 SRP 的核糖体对接到 ER。 \\
Sec61 & ER translocon 核心通道 & 多肽进入 ER 或插入膜。 \\
Sec62 & 辅助转位蛋白 & 参与 post-translational translocation。 \\
Sec63 & 辅助转位蛋白 & 协助 Sec61/BiP 通路。 \\
Sec71 & 辅助转位蛋白 & 参与 post-translational translocation。 \\
Sec72 & 辅助转位蛋白 & 参与 post-translational translocation。 \\
Signal peptidase & 信号肽酶 & 切除进入 ER 的信号肽。 \\
BiP & ER 腔内伴侣蛋白 & 帮助折叠并推动蛋白单向进入 ER。 \\
Calnexin & ER 伴侣蛋白 & 结合糖链，帮助糖蛋白折叠。 \\
Protein disulfide isomerase & 蛋白二硫键异构酶 & 帮助形成/重排二硫键。 \\
Lectin & 凝集素类伴侣 & 识别糖链并帮助 ER 质量控制。 \\
Glucosidase & 葡萄糖苷酶 & 修剪 N-linked 糖链中的葡萄糖。 \\
Glucosyl transferase & 葡萄糖转移酶 & 给未折好蛋白重新加葡萄糖进入 calnexin 循环。 \\
Oligosaccharyl transferase (OST) & 寡糖转移酶 & 把预装配糖链转移到 Asn 上。 \\
IRE1 & ER 应激感受蛋白 & 感知未折叠蛋白并启动 UPR。 \\
E3 ubiquitin ligase & E3泛素连接酶 & 给错误折叠蛋白加 ubiquitin 标记。 \\
Ubiquitin & 泛素 & 标记蛋白用于降解或调控。 \\
N-glycanase & N-聚糖酶 & 去除 ERAD 底物上的糖链。 \\
AAA-ATPase & AAA ATP酶 & 用 ATP 把 ERAD 底物拉出膜。 \\
Proteasome & 蛋白酶体 & 降解被泛素化蛋白。 \\
TMEM41B & 膜蛋白 & 与脂质组织/病毒复制相关。 \\
Torsin & AAA$^+$ ATPase 家族蛋白 & 与核膜/ER 膜系统稳态相关。 \\
WRB & Tail-anchored 蛋白受体组分 & 接收 TRC40 递送的 tail-anchored protein。 \\
CAML & Tail-anchored 蛋白受体组分 & 与 WRB 组成受体复合体。 \\
TRC40 & 递送 ATPase & 把 tail-anchored protein 送到 ER 膜。 \\
Get1 & TA蛋白插入受体组分 & 接收 Get3 递送的 tail-anchored protein。 \\
Get2 & TA蛋白插入受体组分 & 与 Get1 协同插入 tail-anchored protein。 \\
Get3 & ATPase & 递送 tail-anchored protein 到 ER 受体。 \\
SGTA & 捕获蛋白 & 先结合 tail-anchored protein 的疏水尾段。 \\
Bag6 & 分选/质控蛋白 & 参与 tail-anchored protein 分选与质量控制。 \\
Catalase & 过氧化氢酶 & 分解 H$_2$O$_2$。 \\
Urate oxidase & 尿酸氧化酶 & 参与过氧化物酶体氧化代谢。 \\
\end{longtable}

\section*{6. 囊泡运输与膜交通}
\begin{longtable}{L{3.2cm} L{3.8cm} L{8.5cm}}
\toprule
\textbf{名称} & \textbf{中文/类别} & \textbf{功能（极简）} \\
\midrule
\endfirsthead
\toprule
\textbf{名称} & \textbf{中文/类别} & \textbf{功能（极简）} \\
\midrule
\endhead
\bottomrule
\endfoot
Clathrin & 网格蛋白 & 形成 clathrin-coated vesicle 外壳。 \\
Adaptor proteins & 接头蛋白 & 选择 cargo 并连接 cargo receptor 与 coat。 \\
Sar1 & 小GTP酶 & 招募 COPII coat 到 ER 膜。 \\
Sec23 & COPII 内层被覆蛋白 & 与 Sar1 结合并参与 cargo 选择。 \\
Sec24 & COPII 内层被覆蛋白 & 识别 cargo。 \\
COPI coat proteins & COPI 被覆蛋白 & 介导从 Golgi 回 ER 的回收运输。 \\
Rab proteins & Rab 小GTP酶 & 指定囊泡靶向与 docking 位置。 \\
Rab effectors & Rab 效应蛋白 & tether 囊泡到正确靶膜。 \\
GDI & Rab-GDP dissociation inhibitor & 结合 Rab-GDP，使其在胞质中可溶。 \\
SNARE proteins & SNARE 蛋白 & 直接催化膜融合。 \\
v-SNARE & 囊泡 SNARE & 位于囊泡膜上，与 t-SNARE 配对。 \\
t-SNARE & 靶膜 SNARE & 位于靶膜上，与 v-SNARE 配对。 \\
NSF & ATPase & 解离已经形成的 SNARE 复合体。 \\
Complexin & SNARE 调控蛋白 & 稳定 SNARE 复合体并防止过早融合。 \\
BiP (ER retrieval context) & ER 驻留蛋白 & 也是 KDEL 类回收路径 cargo 的例子。 \\
Mannose 6-phosphate receptor & M6P 受体 & 在 TGN 识别 lysosomal hydrolase。 \\
GlcNAc phosphotransferase & GlcNAc 磷酸转移酶 & 给 lysosomal hydrolase 打 M6P 前体标记。 \\
Retromer & 回收被覆/复合体 & 参与从内体回收受体等货物。 \\
\end{longtable}

\section*{7. 术语与方法（来自 notebook，包括图片相关内容）}
\begin{longtable}{L{3.2cm} L{3.8cm} L{8.5cm}}
\toprule
\textbf{名称} & \textbf{中文/类别} & \textbf{功能/含义（极简）} \\
\midrule
\endfirsthead
\toprule
\textbf{名称} & \textbf{中文/类别} & \textbf{功能/含义（极简）} \\
\midrule
\endhead
\bottomrule
\endfoot
Heterokaryon & 异核体 & 两个细胞融合后、同一细胞质中暂时保留两个独立细胞核的状态。 \\
Light microscope & 光学显微镜 & 用可见光成像。 \\
Upright microscope & 正置显微镜 & 物镜在样品上方，常用于常规切片观察。 \\
Inverted microscope & 倒置显微镜 & 物镜在样品下方，常用于细胞培养观察。 \\
Dark-field microscope & 暗场显微镜 & 利用散射光在黑背景上成像。 \\
Phase-contrast microscopy & 相差显微镜 & 把相位差转成明暗差，适合透明活细胞。 \\
Fluorescence microscopy & 荧光显微镜 & 检测特定荧光分子/标记。 \\
Conventional fluorescence microscopy & 常规荧光显微镜 & 普通荧光成像，易受离焦光干扰。 \\
Confocal microscopy & 共聚焦显微镜 & 排除离焦光，获得光学切片。 \\
Super-resolution microscopy & 超分辨显微镜 & 分辨率突破传统光学极限。 \\
Atomic force microscopy & 原子力显微镜 & 用机械探针扫描表面形貌。 \\
Electron microscope & 电子显微镜 & 用电子束成像，分辨率高。 \\
TEM & 透射电子显微镜 & 用穿透样品的电子成像内部结构。 \\
SEM & 扫描电子显微镜 & 扫描样品表面，观察表面形貌。 \\
Cryo-electron microscopy & 冷冻电镜 & 在玻璃态冰中对样品低温成像并重建结构。 \\
Immunogold EM & 免疫金电镜 & 用带金颗粒抗体在电镜下定位蛋白。 \\
Optical section & 光学切片 & 通过成像手段得到样品某一薄层的信息。 \\
Deconvolution & 去卷积 & 计算去除离焦荧光，提高图像清晰度。 \\
FRET & 荧光共振能量转移 & 测分子间近距离相互作用。 \\
FRAP & 光漂白后荧光恢复 & 测膜蛋白/脂质扩散和流动性。 \\
Flow cytometry & 流式细胞术 & 按散射和荧光分析单细胞。 \\
FACS & 荧光激活细胞分选 & 按荧光等特征分选细胞。 \\
In situ hybridization & 原位杂交 & 在组织/细胞原位检测特定核酸。 \\
FISH & 荧光原位杂交 & 用荧光探针定位核酸或染色体位点。 \\
Immunostaining & 免疫染色 & 用抗体检测特定分子。 \\
Direct antibody staining & 直接抗体染色 & 带标记的一抗直接识别目标。 \\
Indirect antibody staining & 间接抗体染色 & 二抗放大信号。 \\
Polyclonal antibody & 多克隆抗体 & 来自多个 B 细胞克隆，识别多个表位。 \\
Monoclonal antibody & 单克隆抗体 & 来自单个 B 细胞克隆，特异性高。 \\
Hybridoma & 杂交瘤 & 稳定产生单克隆抗体的细胞系。 \\
Cell fractionation & 细胞组分分离 & 把细胞分成不同亚细胞组分。 \\
Centrifugation & 离心分离 & 依靠离心力分离颗粒。 \\
Preparative ultracentrifuge & 制备型超速离心 & 用于分离和收集组分。 \\
Analytical ultracentrifuge & 分析型超速离心 & 用于分析沉降性质。 \\
Velocity sedimentation & 速度沉降 & 按沉降速度分离颗粒。 \\
Equilibrium sedimentation & 平衡沉降 & 按浮力密度分离颗粒。 \\
Chromatography & 色谱法 & 利用分配差异分离分子。 \\
Sequential chromatography & 串联色谱 & 组合多种色谱连续纯化。 \\
Ion-exchange chromatography & 离子交换色谱 & 按电荷分离。 \\
Hydrophobic chromatography & 疏水色谱 & 按疏水性分离。 \\
Gel-filtration chromatography & 凝胶过滤/分子筛色谱 & 按大小分离。 \\
Affinity chromatography & 亲和色谱 & 按特异结合能力分离。 \\
Immunoprecipitation & 免疫沉淀 & 用抗体富集目标蛋白。 \\
Co-immunoprecipitation (Co-IP) & 免疫共沉淀 & 检测蛋白-蛋白相互作用。 \\
Tandem affinity purification (TAP) & 串联亲和纯化 & 用双标签高纯度纯化复合体。 \\
SDS-PAGE & SDS 聚丙烯酰胺凝胶电泳 & 按蛋白大小分离。 \\
Two-dimensional gel electrophoresis & 双向凝胶电泳 & 先按电荷再按大小分离蛋白。 \\
Blotting & 印迹法 & 转膜后用探针/抗体检测分子。 \\
Mass spectrometry & 质谱 & 按质荷比分离并鉴定分子。 \\
LC-MS & 液相色谱-质谱联用 & 先色谱后质谱，分析复杂蛋白混合物。 \\
X-ray diffraction & X 射线衍射 & 用衍射图重建蛋白结构。 \\
Recombinant DNA technology & 重组 DNA 技术 & 精确操作 DNA 的一组方法。 \\
Restriction nuclease & 限制性核酸内切酶 & 在特定序列切割 DNA。 \\
DNA ligation & DNA 连接 & 连接 DNA 片段。 \\
DNA cloning & DNA 克隆 & 扩增并保存特定 DNA 片段。 \\
DNA library & DNA 文库 & 某一来源 DNA 片段的集合。 \\
Genomic library & 基因组文库 & 覆盖整个基因组 DNA。 \\
cDNA library & cDNA 文库 & 反映表达过 RNA 的 DNA 文库。 \\
PCR & 聚合酶链式反应 & 指数扩增特定 DNA 片段。 \\
Sanger sequencing & 桑格测序 & 用 ddNTP 终止法读取序列。 \\
Second-generation sequencing & 二代测序 & 高通量并行测序。 \\
Nucleic acid hybridization & 核酸杂交 & 用互补配对检测特定序列。 \\
Microarray & 芯片表达谱 & 同时检测大量基因表达。 \\
RNA-seq & RNA 测序 & 高通量分析转录组。 \\
Genome annotation & 基因组注释 & 标记基因与功能元件。 \\
Genetic screen & 遗传筛选 & 寻找具有特定表型的突变体。 \\
Reverse genetics & 反向遗传学 & 从已知基因出发研究功能。 \\
Knock-out (KO) & 基因敲除 & 去除基因功能。 \\
Knock-in (KI) & 基因敲入 & 定点引入新序列。 \\
Conditional KO/KI/Expr & 条件性遗传操作 & 在特定条件/组织中调控基因。 \\
CRISPR-mediated immunity & CRISPR 免疫 & 原核生物利用 CRISPR 抵御外源核酸。 \\
RNA interference & RNA 干扰 & 用小 RNA 抑制基因表达。 \\
VNTR & 可变串联重复序列 & 个体间高度多态，常用于法医学。 \\
STR & 短串联重复序列 & 常用遗传标记和法医学标记。 \\
Polymorphism & 多态性 & 群体中常见的序列变异。 \\
Haplotype block & 单倍型块 & 常一起遗传的一组多态位点。 \\
Linkage disequilibrium & 连锁不平衡 & 等位位点组合偏离随机。 \\
\end{longtable}

\section*{8. 补充说明}
\begin{itemize}
\item 若你想要更“考试友好”的版本，我可以继续把这张表改成四列：\textbf{名称、定位、功能、常见考法}。
\item 若你想要更“背诵友好”的版本，我也可以再导出一个\textbf{只保留英文名 + 一句话功能}的超精简版。
\end{itemize}

\end{document}
